\documentclass[11pt,a4paper]{article}
% ==============================================================
%  Préambule commun aux fiches -- Advanced Probabilities (M1)
%  Usage, dans chaque fiche :
%     \documentclass[11pt,a4paper]{article}
%     % ==============================================================
%  Préambule commun aux fiches -- Advanced Probabilities (M1)
%  Usage, dans chaque fiche :
%     \documentclass[11pt,a4paper]{article}
%     % ==============================================================
%  Préambule commun aux fiches -- Advanced Probabilities (M1)
%  Usage, dans chaque fiche :
%     \documentclass[11pt,a4paper]{article}
%     \input{../preambule.tex}
%     \fiche{C5}{Vecteurs gaussiens}{Gaussian vectors}
% ==============================================================
\usepackage[T1]{fontenc}
\usepackage[utf8]{inputenc}
\usepackage[french]{babel}
\usepackage{lmodern}
\usepackage[a4paper,margin=1.9cm,top=1.7cm,bottom=1.9cm]{geometry}
\usepackage{amsmath,amssymb,amsthm,mathtools}
\usepackage[inline]{enumitem}
\usepackage{booktabs,array,tabularx}
\usepackage{xcolor}
\usepackage[most]{tcolorbox}
\usepackage{fancyhdr}
\usepackage{hyperref}
\hypersetup{colorlinks=true,linkcolor=blue!50!black,urlcolor=blue!50!black}

\setlist{itemsep=1pt,topsep=3pt,parsep=0pt}
\setlength{\parindent}{0pt}
\setlength{\parskip}{3pt}
\emergencystretch=3em

% ---------- Identification de la fiche -------------------------
\newcommand{\ficheid}{}
\newcommand{\fichetitre}{}
\newcommand{\fiche}[3]{%
  \renewcommand{\ficheid}{#1}\renewcommand{\fichetitre}{#2}%
  \pagestyle{fancy}\fancyhf{}%
  \fancyhead[L]{\small\textsf{Advanced Probabilities -- Fiche #1}}%
  \fancyhead[R]{\small\textsf{#2}}%
  \fancyfoot[C]{\small\thepage}%
  \begin{center}
    {\LARGE\bfseries Fiche #1 -- #2}\par\smallskip
    {\large\itshape (#3)}\par\smallskip
    {\small Cours de V.\,M.~Lambert, M1 -- semestre 1}
  \end{center}\vspace{-2pt}\hrule\vspace{6pt}}

% ---------- Environnements numérotés (un seul compteur) --------
\theoremstyle{definition}
\newtheorem{definition}{Définition}
\renewcommand{\thedefinition}{\ficheid.\arabic{definition}}
\newtheorem{theoreme}[definition]{Théorème}
\newtheorem{proposition}[definition]{Proposition}
\newtheorem{lemme}[definition]{Lemme}
\newtheorem{corollaire}[definition]{Corollaire}
\newtheorem{exemple}[definition]{Exemple}
\theoremstyle{remark}
\newtheorem{remarque}[definition]{Remarque}

% Résultat hors cours (standard, signalé comme tel)
\theoremstyle{definition}
\newtheorem{horscours}[definition]{Hors cours}

% Idée de preuve
\newenvironment{idee}{\par\small\textit{Idée de preuve.}~}{\par}

% ---------- Encadrés -----------------------------------------
\tcbset{fichebox/.style={enhanced,breakable,boxrule=0.6pt,arc=2pt,
  left=5pt,right=5pt,top=3pt,bottom=3pt,fonttitle=\bfseries}}
\newtcolorbox{piege}[1][]{fichebox,colback=red!4,colframe=red!60!black,
  title={Pièges classiques},#1}
\newtcolorbox{liens}[1][]{fichebox,colback=blue!3,colframe=blue!50!black,
  title={Liens entre notions},#1}
\newtcolorbox{methode}[1][Méthode]{fichebox,colback=green!4,colframe=green!40!black,title={#1}}
\newtcolorbox{alerte}[1][]{fichebox,colback=orange!6,colframe=orange!70!black,
  title={Attention à l'énoncé},#1}

% Exercice de TD (fiches de rappels)
\newcounter{exo}
\newenvironment{exo}[1][]{\refstepcounter{exo}\par\medskip
  \noindent{\large\bfseries Exercice~\theexo}\ifx\relax#1\relax\else~{\large\bfseries -- #1}\fi\par\nopagebreak\smallskip}{\par}
\newcommand{\outils}{\textbf{Outils du cours.}~}
\newcommand{\meth}{\textbf{Méthode / astuce.}~}
\newcommand{\indic}{\textbf{Indication.}~}
\newcommand{\rep}{\textbf{Réponse.}~}

% ---------- Notations du cours ---------------------------------
\newcommand{\R}{\mathbb{R}}
\newcommand{\N}{\mathbb{N}}
\newcommand{\Q}{\mathbb{Q}}
\newcommand{\C}{\mathbb{C}}
\newcommand{\PP}{\mathbb{P}}
\newcommand{\QQ}{\mathbb{Q}}
\newcommand{\E}{\mathbb{E}}
\newcommand{\F}{\mathcal{F}}
\newcommand{\G}{\mathcal{G}}
\newcommand{\Hc}{\mathcal{H}}
\newcommand{\Tc}{\mathcal{T}}
\newcommand{\Bc}{\mathcal{B}}
\newcommand{\Cc}{\mathcal{C}}
\newcommand{\Ec}{\mathcal{E}}
\newcommand{\Pc}{\mathcal{P}}
\newcommand{\Nc}{\mathcal{N}}
\newcommand{\LN}{\mathcal{LN}}
\newcommand{\Expo}{\mathcal{E}}
\newcommand{\Poi}{\mathcal{P}}
\newcommand{\Ff}{\mathbb{F}}
\newcommand{\1}{\mathbf{1}}
\newcommand{\dd}{\,\mathrm{d}}
\DeclareMathOperator{\Var}{\mathbb{V}}
\DeclareMathOperator{\Cov}{Cov}
\DeclareMathOperator{\card}{card}
\DeclareMathOperator{\diag}{diag}
\newcommand{\indep}{\perp\!\!\!\perp}
\newcommand{\ps}{\ \PP\text{-p.s.}}
\newcommand{\tendloi}{\xrightarrow[n\to+\infty]{\text{loi}}}
\newcommand{\tendps}{\xrightarrow[n\to+\infty]{\text{p.s.}}}
\newcommand{\tendproba}{\xrightarrow[n\to+\infty]{\PP}}
\newcommand{\tendLp}{\xrightarrow[n\to+\infty]{L^p}}
\newcommand{\tend}{\xrightarrow[n\to+\infty]{}}

%     \fiche{C5}{Vecteurs gaussiens}{Gaussian vectors}
% ==============================================================
\usepackage[T1]{fontenc}
\usepackage[utf8]{inputenc}
\usepackage[french]{babel}
\usepackage{lmodern}
\usepackage[a4paper,margin=1.9cm,top=1.7cm,bottom=1.9cm]{geometry}
\usepackage{amsmath,amssymb,amsthm,mathtools}
\usepackage[inline]{enumitem}
\usepackage{booktabs,array,tabularx}
\usepackage{xcolor}
\usepackage[most]{tcolorbox}
\usepackage{fancyhdr}
\usepackage{hyperref}
\hypersetup{colorlinks=true,linkcolor=blue!50!black,urlcolor=blue!50!black}

\setlist{itemsep=1pt,topsep=3pt,parsep=0pt}
\setlength{\parindent}{0pt}
\setlength{\parskip}{3pt}
\emergencystretch=3em

% ---------- Identification de la fiche -------------------------
\newcommand{\ficheid}{}
\newcommand{\fichetitre}{}
\newcommand{\fiche}[3]{%
  \renewcommand{\ficheid}{#1}\renewcommand{\fichetitre}{#2}%
  \pagestyle{fancy}\fancyhf{}%
  \fancyhead[L]{\small\textsf{Advanced Probabilities -- Fiche #1}}%
  \fancyhead[R]{\small\textsf{#2}}%
  \fancyfoot[C]{\small\thepage}%
  \begin{center}
    {\LARGE\bfseries Fiche #1 -- #2}\par\smallskip
    {\large\itshape (#3)}\par\smallskip
    {\small Cours de V.\,M.~Lambert, M1 -- semestre 1}
  \end{center}\vspace{-2pt}\hrule\vspace{6pt}}

% ---------- Environnements numérotés (un seul compteur) --------
\theoremstyle{definition}
\newtheorem{definition}{Définition}
\renewcommand{\thedefinition}{\ficheid.\arabic{definition}}
\newtheorem{theoreme}[definition]{Théorème}
\newtheorem{proposition}[definition]{Proposition}
\newtheorem{lemme}[definition]{Lemme}
\newtheorem{corollaire}[definition]{Corollaire}
\newtheorem{exemple}[definition]{Exemple}
\theoremstyle{remark}
\newtheorem{remarque}[definition]{Remarque}

% Résultat hors cours (standard, signalé comme tel)
\theoremstyle{definition}
\newtheorem{horscours}[definition]{Hors cours}

% Idée de preuve
\newenvironment{idee}{\par\small\textit{Idée de preuve.}~}{\par}

% ---------- Encadrés -----------------------------------------
\tcbset{fichebox/.style={enhanced,breakable,boxrule=0.6pt,arc=2pt,
  left=5pt,right=5pt,top=3pt,bottom=3pt,fonttitle=\bfseries}}
\newtcolorbox{piege}[1][]{fichebox,colback=red!4,colframe=red!60!black,
  title={Pièges classiques},#1}
\newtcolorbox{liens}[1][]{fichebox,colback=blue!3,colframe=blue!50!black,
  title={Liens entre notions},#1}
\newtcolorbox{methode}[1][Méthode]{fichebox,colback=green!4,colframe=green!40!black,title={#1}}
\newtcolorbox{alerte}[1][]{fichebox,colback=orange!6,colframe=orange!70!black,
  title={Attention à l'énoncé},#1}

% Exercice de TD (fiches de rappels)
\newcounter{exo}
\newenvironment{exo}[1][]{\refstepcounter{exo}\par\medskip
  \noindent{\large\bfseries Exercice~\theexo}\ifx\relax#1\relax\else~{\large\bfseries -- #1}\fi\par\nopagebreak\smallskip}{\par}
\newcommand{\outils}{\textbf{Outils du cours.}~}
\newcommand{\meth}{\textbf{Méthode / astuce.}~}
\newcommand{\indic}{\textbf{Indication.}~}
\newcommand{\rep}{\textbf{Réponse.}~}

% ---------- Notations du cours ---------------------------------
\newcommand{\R}{\mathbb{R}}
\newcommand{\N}{\mathbb{N}}
\newcommand{\Q}{\mathbb{Q}}
\newcommand{\C}{\mathbb{C}}
\newcommand{\PP}{\mathbb{P}}
\newcommand{\QQ}{\mathbb{Q}}
\newcommand{\E}{\mathbb{E}}
\newcommand{\F}{\mathcal{F}}
\newcommand{\G}{\mathcal{G}}
\newcommand{\Hc}{\mathcal{H}}
\newcommand{\Tc}{\mathcal{T}}
\newcommand{\Bc}{\mathcal{B}}
\newcommand{\Cc}{\mathcal{C}}
\newcommand{\Ec}{\mathcal{E}}
\newcommand{\Pc}{\mathcal{P}}
\newcommand{\Nc}{\mathcal{N}}
\newcommand{\LN}{\mathcal{LN}}
\newcommand{\Expo}{\mathcal{E}}
\newcommand{\Poi}{\mathcal{P}}
\newcommand{\Ff}{\mathbb{F}}
\newcommand{\1}{\mathbf{1}}
\newcommand{\dd}{\,\mathrm{d}}
\DeclareMathOperator{\Var}{\mathbb{V}}
\DeclareMathOperator{\Cov}{Cov}
\DeclareMathOperator{\card}{card}
\DeclareMathOperator{\diag}{diag}
\newcommand{\indep}{\perp\!\!\!\perp}
\newcommand{\ps}{\ \PP\text{-p.s.}}
\newcommand{\tendloi}{\xrightarrow[n\to+\infty]{\text{loi}}}
\newcommand{\tendps}{\xrightarrow[n\to+\infty]{\text{p.s.}}}
\newcommand{\tendproba}{\xrightarrow[n\to+\infty]{\PP}}
\newcommand{\tendLp}{\xrightarrow[n\to+\infty]{L^p}}
\newcommand{\tend}{\xrightarrow[n\to+\infty]{}}

%     \fiche{C5}{Vecteurs gaussiens}{Gaussian vectors}
% ==============================================================
\usepackage[T1]{fontenc}
\usepackage[utf8]{inputenc}
\usepackage[french]{babel}
\usepackage{lmodern}
\usepackage[a4paper,margin=1.9cm,top=1.7cm,bottom=1.9cm]{geometry}
\usepackage{amsmath,amssymb,amsthm,mathtools}
\usepackage[inline]{enumitem}
\usepackage{booktabs,array,tabularx}
\usepackage{xcolor}
\usepackage[most]{tcolorbox}
\usepackage{fancyhdr}
\usepackage{hyperref}
\hypersetup{colorlinks=true,linkcolor=blue!50!black,urlcolor=blue!50!black}

\setlist{itemsep=1pt,topsep=3pt,parsep=0pt}
\setlength{\parindent}{0pt}
\setlength{\parskip}{3pt}
\emergencystretch=3em

% ---------- Identification de la fiche -------------------------
\newcommand{\ficheid}{}
\newcommand{\fichetitre}{}
\newcommand{\fiche}[3]{%
  \renewcommand{\ficheid}{#1}\renewcommand{\fichetitre}{#2}%
  \pagestyle{fancy}\fancyhf{}%
  \fancyhead[L]{\small\textsf{Advanced Probabilities -- Fiche #1}}%
  \fancyhead[R]{\small\textsf{#2}}%
  \fancyfoot[C]{\small\thepage}%
  \begin{center}
    {\LARGE\bfseries Fiche #1 -- #2}\par\smallskip
    {\large\itshape (#3)}\par\smallskip
    {\small Cours de V.\,M.~Lambert, M1 -- semestre 1}
  \end{center}\vspace{-2pt}\hrule\vspace{6pt}}

% ---------- Environnements numérotés (un seul compteur) --------
\theoremstyle{definition}
\newtheorem{definition}{Définition}
\renewcommand{\thedefinition}{\ficheid.\arabic{definition}}
\newtheorem{theoreme}[definition]{Théorème}
\newtheorem{proposition}[definition]{Proposition}
\newtheorem{lemme}[definition]{Lemme}
\newtheorem{corollaire}[definition]{Corollaire}
\newtheorem{exemple}[definition]{Exemple}
\theoremstyle{remark}
\newtheorem{remarque}[definition]{Remarque}

% Résultat hors cours (standard, signalé comme tel)
\theoremstyle{definition}
\newtheorem{horscours}[definition]{Hors cours}

% Idée de preuve
\newenvironment{idee}{\par\small\textit{Idée de preuve.}~}{\par}

% ---------- Encadrés -----------------------------------------
\tcbset{fichebox/.style={enhanced,breakable,boxrule=0.6pt,arc=2pt,
  left=5pt,right=5pt,top=3pt,bottom=3pt,fonttitle=\bfseries}}
\newtcolorbox{piege}[1][]{fichebox,colback=red!4,colframe=red!60!black,
  title={Pièges classiques},#1}
\newtcolorbox{liens}[1][]{fichebox,colback=blue!3,colframe=blue!50!black,
  title={Liens entre notions},#1}
\newtcolorbox{methode}[1][Méthode]{fichebox,colback=green!4,colframe=green!40!black,title={#1}}
\newtcolorbox{alerte}[1][]{fichebox,colback=orange!6,colframe=orange!70!black,
  title={Attention à l'énoncé},#1}

% Exercice de TD (fiches de rappels)
\newcounter{exo}
\newenvironment{exo}[1][]{\refstepcounter{exo}\par\medskip
  \noindent{\large\bfseries Exercice~\theexo}\ifx\relax#1\relax\else~{\large\bfseries -- #1}\fi\par\nopagebreak\smallskip}{\par}
\newcommand{\outils}{\textbf{Outils du cours.}~}
\newcommand{\meth}{\textbf{Méthode / astuce.}~}
\newcommand{\indic}{\textbf{Indication.}~}
\newcommand{\rep}{\textbf{Réponse.}~}

% ---------- Notations du cours ---------------------------------
\newcommand{\R}{\mathbb{R}}
\newcommand{\N}{\mathbb{N}}
\newcommand{\Q}{\mathbb{Q}}
\newcommand{\C}{\mathbb{C}}
\newcommand{\PP}{\mathbb{P}}
\newcommand{\QQ}{\mathbb{Q}}
\newcommand{\E}{\mathbb{E}}
\newcommand{\F}{\mathcal{F}}
\newcommand{\G}{\mathcal{G}}
\newcommand{\Hc}{\mathcal{H}}
\newcommand{\Tc}{\mathcal{T}}
\newcommand{\Bc}{\mathcal{B}}
\newcommand{\Cc}{\mathcal{C}}
\newcommand{\Ec}{\mathcal{E}}
\newcommand{\Pc}{\mathcal{P}}
\newcommand{\Nc}{\mathcal{N}}
\newcommand{\LN}{\mathcal{LN}}
\newcommand{\Expo}{\mathcal{E}}
\newcommand{\Poi}{\mathcal{P}}
\newcommand{\Ff}{\mathbb{F}}
\newcommand{\1}{\mathbf{1}}
\newcommand{\dd}{\,\mathrm{d}}
\DeclareMathOperator{\Var}{\mathbb{V}}
\DeclareMathOperator{\Cov}{Cov}
\DeclareMathOperator{\card}{card}
\DeclareMathOperator{\diag}{diag}
\newcommand{\indep}{\perp\!\!\!\perp}
\newcommand{\ps}{\ \PP\text{-p.s.}}
\newcommand{\tendloi}{\xrightarrow[n\to+\infty]{\text{loi}}}
\newcommand{\tendps}{\xrightarrow[n\to+\infty]{\text{p.s.}}}
\newcommand{\tendproba}{\xrightarrow[n\to+\infty]{\PP}}
\newcommand{\tendLp}{\xrightarrow[n\to+\infty]{L^p}}
\newcommand{\tend}{\xrightarrow[n\to+\infty]{}}

\begin{document}
\fiche{C7}{Changement de loi}{Change of law}

\textit{Source : slides, chapitre 8 (Déf.~57, Thm~16, 17) ; poly « Gaussian vector » (Déf.~22, Thm~23, 24).
Motivation : en finance coexistent la probabilité « historique » (\emph{physical}) et la probabilité « risque-neutre »
(\emph{pricing / martingale measure}) ; on veut passer de l'une à l'autre en gardant les mêmes événements impossibles et certains.}

\section*{1. Comparaison de deux probabilités}

\begin{definition}[absolue continuité, équivalence, singularité]\label{C7:def}
Soient $\PP,\QQ$ deux probabilités sur un même espace mesurable $(\Omega,\F)$.
\begin{itemize}
  \item $\QQ$ est \emph{absolument continue} par rapport à $\PP$, noté $\QQ\ll\PP$, si
  $\forall A\in\F$, $\PP(A)=0\Rightarrow\QQ(A)=0$. Les événements impossibles (resp.\ certains) sous $\PP$ le restent sous $\QQ$.
  \item $\PP$ et $\QQ$ sont \emph{équivalentes}, $\PP\sim\QQ$, si $\PP\ll\QQ$ et $\QQ\ll\PP$ : mêmes événements impossibles (et certains).
  \item $\PP$ et $\QQ$ sont \emph{étrangères} (\emph{singular}), $\PP\perp\QQ$, s'il existe $N\in\F$ avec $\PP(N)=0$ et $\QQ(N^c)=0$.
\end{itemize}
\end{definition}

\begin{exemple}[slides]
$\lambda\perp\delta_0$ sur $(\R,\Bc(\R))$ : $N=\{0\}$ vérifie $\lambda(\{0\})=0$ et $\delta_0(\{0\}^c)=0$.
(Les définitions s'étendent aux mesures quelconques ; $\lambda$ n'est pas une probabilité.)
\end{exemple}

\section*{2. Théorème de Radon--Nikodym}

\begin{proposition}\label{C7:densite-ac}
Si $\QQ$ a une densité $f$ (mesurable, positive) par rapport à $\PP$ (i.e.\ $\QQ(A)=\int_A f\dd\PP$), alors $\QQ\ll\PP$ :
si $\PP(A)=0$, $\QQ(A)=\int\1_Af\dd\PP=0$.
\end{proposition}

\begin{theoreme}[Radon--Nikodym]\label{C7:RN}
Pour $\QQ,\PP$ probabilités sur $(\Omega,\F)$, sont équivalents :
\begin{enumerate}
  \item $\QQ\ll\PP$ ;
  \item il existe $f\in L^1_{\R_+}(\Omega,\F,\PP)$, \textbf{unique $\PP$-p.s.}, telle que $\QQ=f\cdot\PP$, i.e.\
  $\QQ(A)=\int\1_Af\dd\PP=\E_\PP(\1_Af)$ pour tout $A\in\F$.
\end{enumerate}
$f$ est la \emph{dérivée de Radon--Nikodym}, notée $f=\frac{\dd\QQ}{\dd\PP}$. Cas discret : $\QQ(\{\omega\})=f(\omega)\PP(\{\omega\})$.
\end{theoreme}

\begin{corollaire}[standard]\label{C7:EQ}
Si $\QQ=f\cdot\PP$ : $\E_\PP(f)=\QQ(\Omega)=1$ ; pour $X\ge0$ (ou $Xf\in L^1(\PP)$) : $\E_\QQ(X)=\E_\PP(Xf)$.
\emph{Hors cours :} $\QQ\sim\PP\iff f>0$ $\PP$-p.s., et alors $\frac{\dd\PP}{\dd\QQ}=\frac1f$.
\end{corollaire}
\begin{idee}
Vrai pour $X=\1_A$ (définition), puis fonctions étagées et limites croissantes (même schéma que le transfert, C2).
\end{idee}

\section*{3. Formule de Bayes conditionnelle}

\begin{theoreme}[Bayes, \emph{Bayes formula}]\label{C7:bayes}
Soit $\G\subset\F$ une sous-tribu et $L:\Omega\to\R_+^*$ une v.a.\ telle que $\E_\PP(L)=1$. On pose $\QQ(A)=\E_\PP(\1_AL)$ :
$\QQ$ est une probabilité de densité $L=\frac{\dd\QQ}{\dd\PP}$ ($\QQ=L\cdot\PP$). Pour toute v.a.\ $X$ réelle (ou vecteur de $\R^d$)
$\QQ$-intégrable,
\[\E_\QQ(X\mid\G)=\frac{\E_\PP(XL\mid\G)}{\E_\PP(L\mid\G)}\qquad\PP\text{-p.s.}\]
\end{theoreme}
\begin{idee}
Soit $Z$ le membre de droite ($\G$-mesurable ; $\E_\PP(L\mid\G)>0$ car $L>0$). Pour $G\in\G$ :
$\E_\QQ(\1_GZ)=\E_\PP(\1_GZL)=\E_\PP(\1_GZ\,\E_\PP(L\mid\G))=\E_\PP(\1_G\E_\PP(XL\mid\G))=\E_\PP(\1_GXL)=\E_\QQ(\1_GX)$
(tour et « sortir ce qui est connu », C6).
\end{idee}

\begin{liens}
\begin{itemize}
  \item Densité $\Rightarrow$ absolue continuité (Prop.~\ref{C7:densite-ac}) ; Radon--Nikodym donne la réciproque.
  \item $\QQ\ll\PP$ $\not\Rightarrow$ $\PP\ll\QQ$ : sur $\{U,D\}$, $\PP(U)=\PP(D)=\frac12$ et $\QQ=\delta_U$. On a $\QQ\ll\PP$
  (le seul ensemble $\PP$-négligeable est $\emptyset$), mais $\QQ(\{D\})=0<\PP(\{D\})$.
  \item Équivalentes $\neq$ égales : mêmes ensembles négligeables, probabilités différentes.
  \item Le théorème de Radon--Nikodym donne l'existence de l'espérance conditionnelle dans $L^1$ (C6).
\end{itemize}
\end{liens}

\begin{piege}
\begin{itemize}
  \item $\ll$ n'est pas symétrique ; préciser « par rapport à quelle probabilité ».
  \item Dans Bayes, diviser par $\E_\PP(L\mid\G)$ (une v.a.), \emph{pas} par $\E_\PP(L)=1$.
  \item Une espérance change avec la probabilité : écrire $\E_\PP$ ou $\E_\QQ$, jamais $\E$ tout court quand deux probabilités coexistent.
  \item La propriété de martingale et l'indépendance dépendent de la probabilité (C3, C8).
  \item La densité $\frac{\dd\QQ}{\dd\PP}$ est définie $\PP$-p.s.\ seulement.
\end{itemize}
\end{piege}

\section*{Exemples types}

\begin{exemple}[deux états, slides]\label{C7:ex-UD}
$\Omega=\{U,D\}$, $\PP(\{U\})=p=1-\PP(\{D\})$, $0<p<1$, et $\QQ$ « équitable » : $\QQ(\{U\})=\QQ(\{D\})=\frac12$.
$f(U)=\frac{\QQ(\{U\})}{\PP(\{U\})}=\frac1{2p}$, $f(D)=\frac1{2(1-p)}$, et $\QQ(\omega)=f(\omega)\PP(\omega)$.
Comme $\PP(\{U\}),\PP(\{D\})>0$ (et de même pour $\QQ$), $\PP\sim\QQ$.
\end{exemple}

\begin{exemple}[Bayes sur l'arbre à deux périodes]\label{C7:ex-arbre}
Modèle C1 sous $\PP$ (paramètre $0<p<1$, pour que $L$ soit défini) ; sous $\QQ$ chaque chemin a probabilité $\frac14$, donc $L(\omega)=\frac{1/4}{\PP(\{\omega\})}$.
Sur $A_U$ : $\E_\PP(S_2L\1_{A_U})=\frac14s_0u^2+\frac14s_0ud$ et $\E_\PP(L\1_{A_U})=\QQ(A_U)=\frac12$, d'où
$\E_\QQ(S_2\mid\F_1)=\frac{u+d}{2}\,s_0u$ sur $A_U$ et, de même, $\frac{u+d}2\,s_0d$ sur $A_D$ : $\E_\QQ(S_2\mid\F_1)=\frac{u+d}2S_1$.
C'est la formule de l'exemple C6 avec $p$ remplacé par $\frac12$.
\end{exemple}

\begin{exemple}[translation gaussienne, hors cours]\label{C7:ex-gauss}
Sous $\PP$, $X\sim\Nc(0,1)$ ; $L=e^{\theta X-\theta^2/2}>0$ vérifie $\E_\PP L=1$. Pour $h\ge0$ :
$\E_\QQ(h(X))=\int h(x)e^{\theta x-\frac{\theta^2}2}\frac{e^{-x^2/2}}{\sqrt{2\pi}}\dd x=\int h(x)\frac{e^{-(x-\theta)^2/2}}{\sqrt{2\pi}}\dd x$ :
sous $\QQ$, $X\sim\Nc(\theta,1)$. On a changé la moyenne sans changer les événements négligeables.
\end{exemple}

\end{document}
